\section{Manipolazione Avanzata di Array e Variadicità (Lez. 17)}

\subsection{Destructuring (Spacchettare gli Array)}
Possiamo estrarre valori da un array e metterli in variabili "spacchettandoli" direttamente con le parentesi quadre.

\begin{concept}{Visualizzazione: Destructuring}
\begin{center}
\begin{tikzpicture}[
    node distance=1cm,
    varbox/.style={draw, rectangle, fill=jsYellow!30, minimum size=0.8cm, font=\ttfamily\bfseries},
    valbox/.style={draw, rectangle, fill=white, minimum size=0.8cm, font=\ttfamily},
    arrow/.style={->, thick, color=violet}
]
    % Variabili
    \node (let) {let [};
    \node[varbox, right=0.1cm of let] (x) {x};
    \node[right=0.1cm of x] (comma) {,};
    \node[varbox, right=0.1cm of comma] (y) {y};
    \node[right=0.1cm of y] (eq) {] = [};
    
    % Valori
    \node[valbox, right=0.1cm of eq] (v1) {10};
    \node[right=0.1cm of v1] (comma2) {,};
    \node[valbox, right=0.1cm of comma2] (v2) {20};
    \node[right=0.1cm of v2] (end) {];};

    % Frecce Mappatura
    \draw[arrow] (v1.north) to[out=90, in=90] (x.north);
    \draw[arrow] (v2.south) to[out=-90, in=-90] (y.south);
\end{tikzpicture}
\end{center}
\end{concept}

\begin{jsbox}
let coordinate = [10, 20];
let [x, y] = coordinate;
console.log(x); // 10
console.log(y); // 20
\end{jsbox}

\subsection{Rest Operator (...)}
I tre puntini \texttt{...} servono a dire "tutto il resto". È utilissimo per prendere la testa e la coda di una lista.

\begin{concept}{Visualizzazione: Rest Operator}
\begin{center}
\begin{tikzpicture}[
    valbox/.style={draw, rectangle, fill=white, minimum size=0.8cm},
    headbox/.style={draw, rectangle, fill=codeGreen!30, minimum size=0.8cm},
    tailbox/.style={draw, rectangle, fill=codePink!30, minimum size=0.8cm}
]
    \node[headbox] (n1) at (0,0) {1};
    \node[tailbox] (n2) at (1,0) {2};
    \node[tailbox] (n3) at (2,0) {3};
    \node[tailbox] (n4) at (3,0) {4};
    
    \node[below=0.5cm of n1, font=\bfseries\color{codeGreen}] {primo};
    
    \draw[decorate, decoration={brace, amplitude=10pt, mirror}, thick] (0.5, -0.2) -- (3.5, -0.2);
    \node[below=0.8cm of n2, xshift=0.5cm, font=\bfseries\color{codePink}] {...altri};
\end{tikzpicture}
\end{center}
\end{concept}

\begin{jsbox}
let numeri = [1, 2, 3, 4, 5];
let [primo, ...altri] = numeri;

console.log(primo); // 1
console.log(altri); // [2, 3, 4, 5]
\end{jsbox}

\subsection{Funzioni Variadiche (Argomenti Variabili)}
Ecco un esempio di funzione che usa il Rest Operator. Usiamo il nuovo **Box Viola** per le funzioni speciali.

\begin{funcbox}[title=Funzione Math Universale]
function matematica(operatore, ...numeri) {
    // 'numeri' diventa un ARRAY con tutti gli extra
    // Esempio: matematica("+", 1, 2, 3) --> numeri = [1, 2, 3]
    
    if (operatore === "+") {
        return numeri.reduce((tot, n) => tot + n, 0);
    }
    if (operatore === "*") {
        return numeri.reduce((tot, n) => tot * n, 1);
    }
}
\end{funcbox}

\begin{jsbox}
console.log(matematica("+", 1, 2, 3)); // 6
console.log(matematica("*", 2, 2, 2, 2)); // 16
\end{jsbox}
